% GENERATED FILE. Edit canonical lessons, then run npm run generate:books.
% canonical-source-hash: fnv1a64:6bbb6e19fd0972af
% canonical-lessons: IT-C374-imprevisto, IT-C374-prevedibile, IT-C374-imprevedibile, IT-C374-alternativa, IT-C374-adattarsi

\chapter{Gli imprevisti}
\label{ch:it-gli-imprevisti}
\hlchaptermodality{374}

\begin{chapteropening}
\textbf{By the end of this chapter:} \emph{I can talk about the unexpected: something foreseeable or not, an alternative and adapting to it.}

Complete the last lesson of chapter 374: I can talk about the unexpected: something foreseeable or not, an alternative and adapting to it.
\end{chapteropening}

\section[l'imprevisto]{l'imprevisto — the unexpected problem, the hitch}

\label{lesson:IT-C374-imprevisto}



\begin{quote}
Before the new one: say the Italian for or else, then the Italian for possible.
\end{quote}

\subsection*{You'll want to know: l'imprevisto}
\textbf{l'imprevisto} — \textquotedblleft{}the unexpected problem, the hitch\textquotedblright{}.

\begin{quote}
Scusa il ritardo, ho avuto un imprevisto.
\end{quote}

\begin{quote}
\emph{Sorry I am late, something came up.}
\end{quote}

As an adjective: unexpected. \emph{Salvo imprevisti} — barring the unexpected.

\begin{cousinweb}
\emph{imprevisto} is built from \emph{in-} \textquotedblleft{}not\textquotedblright{} and \emph{previsto} \textquotedblleft{}foreseen\textquotedblright{}, past participle of \emph{prevedere}, from Latin \emph{praevidēre} \textquotedblleft{}to foresee\textquotedblright{} (Lewis \& Short), on \emph{vidēre} \textquotedblleft{}to see\textquotedblright{}.
\end{cousinweb}

\begin{grammarlens}[title={What you've built}]
One more word for telling, explaining and describing at B1.
\end{grammarlens}

\subsection*{Your turn}
\begin{itemize}
\raggedright
  \item \emph{Say it:} \emph{l'imprevisto}
  \item \emph{Say it:} \emph{l'imprevisto}, once more
  \item \emph{Say it:} a sentence of your own with \emph{l'imprevisto}
  \item \emph{Recall:} say the Italian for or else, then the Italian for possible, then say \emph{l'imprevisto} again
\end{itemize}

\begin{tcolorbox}[breakable,colback=teal!4,colframe=teal!35!black,title={Before you move on}]
What does l'imprevisto mean? (\textquotedblleft{}The unexpected problem, the hitch\textquotedblright{}.) Say it once more.
\end{tcolorbox}
\section[prevedibile]{prevedibile — predictable, foreseeable}

\label{lesson:IT-C374-prevedibile}



\begin{quote}
Before the new one: say the Italian for possible, then the Italian for the unexpected problem.
\end{quote}

\subsection*{You'll want to know: prevedibile}
\textbf{prevedibile} — \textquotedblleft{}predictable, foreseeable\textquotedblright{}.

\begin{quote}
Il finale del film era prevedibile.
\end{quote}

\begin{quote}
\emph{The ending of the film was predictable.}
\end{quote}

Built on \emph{prevedere} \textquotedblleft{}to foresee\textquotedblright{}.

\begin{cousinweb}
\emph{prevedibile} is built on \emph{prevedere} \textquotedblleft{}to foresee\textquotedblright{}, from Latin \emph{praevidēre} \textquotedblleft{}to see beforehand\textquotedblright{} (Lewis \& Short), on \emph{vidēre} \textquotedblleft{}to see\textquotedblright{}.
\end{cousinweb}

\begin{grammarlens}[title={What you've built}]
One more word for telling, explaining and describing at B1.
\end{grammarlens}

\subsection*{Your turn}
\begin{itemize}
\raggedright
  \item \emph{Say it:} \emph{prevedibile}
  \item \emph{Say it:} \emph{prevedibile}, once more
  \item \emph{Say it:} a sentence of your own with \emph{prevedibile}
  \item \emph{Recall:} say the Italian for possible, then the Italian for the unexpected problem, then say \emph{prevedibile} again
\end{itemize}

\begin{tcolorbox}[breakable,colback=teal!4,colframe=teal!35!black,title={Before you move on}]
What does prevedibile mean? (\textquotedblleft{}Predictable, foreseeable\textquotedblright{}.) Say it once more.
\end{tcolorbox}
\section[imprevedibile]{imprevedibile — unpredictable}

\label{lesson:IT-C374-imprevedibile}



\begin{quote}
Before the new one: say the Italian for the unexpected problem, then the Italian for predictable.
\end{quote}

\subsection*{You'll want to know: imprevedibile}
\textbf{imprevedibile} — \textquotedblleft{}unpredictable\textquotedblright{}.

\begin{quote}
Il tempo in montagna è imprevedibile.
\end{quote}

\begin{quote}
\emph{The weather in the mountains is unpredictable.}
\end{quote}

Of a person too: \emph{un carattere imprevedibile}.

\begin{cousinweb}
\emph{imprevedibile} is built from \emph{in-} \textquotedblleft{}not\textquotedblright{} and \emph{prevedibile}, on \emph{prevedere}, from Latin \emph{praevidēre} \textquotedblleft{}to foresee\textquotedblright{} (Lewis \& Short), on \emph{vidēre} \textquotedblleft{}to see\textquotedblright{}.
\end{cousinweb}

\begin{grammarlens}[title={What you've built}]
One more word for telling, explaining and describing at B1.
\end{grammarlens}

\subsection*{Your turn}
\begin{itemize}
\raggedright
  \item \emph{Say it:} \emph{imprevedibile}
  \item \emph{Say it:} \emph{imprevedibile}, once more
  \item \emph{Say it:} a sentence of your own with \emph{imprevedibile}
  \item \emph{Recall:} say the Italian for the unexpected problem, then the Italian for predictable, then say \emph{imprevedibile} again
\end{itemize}

\begin{tcolorbox}[breakable,colback=teal!4,colframe=teal!35!black,title={Before you move on}]
What does imprevedibile mean? (\textquotedblleft{}Unpredictable\textquotedblright{}.) Say it once more.
\end{tcolorbox}
\section[l'alternativa]{l'alternativa — the alternative, the other option}

\label{lesson:IT-C374-alternativa}



\begin{quote}
Before the new one: say the Italian for predictable, then the Italian for unpredictable.
\end{quote}

\subsection*{You'll want to know: l'alternativa}
\textbf{l'alternativa} — \textquotedblleft{}the alternative, the other option\textquotedblright{}.

\begin{quote}
Se il ristorante è pieno, qual è l'alternativa?
\end{quote}

\begin{quote}
\emph{If the restaurant is full, what is the alternative?}
\end{quote}

\emph{Non c'è alternativa} — there is no other way.

\begin{cousinweb}
\emph{alternativa} is a learned borrowing from Medieval Latin \emph{alternātīvus}, from \emph{alternāre} \textquotedblleft{}to do by turns\textquotedblright{} (Lewis \& Short), built on \emph{alter} \textquotedblleft{}the other of two\textquotedblright{}.
\end{cousinweb}

\begin{grammarlens}[title={What you've built}]
One more word for telling, explaining and describing at B1.
\end{grammarlens}

\subsection*{Your turn}
\begin{itemize}
\raggedright
  \item \emph{Say it:} \emph{l'alternativa}
  \item \emph{Say it:} \emph{l'alternativa}, once more
  \item \emph{Say it:} a sentence of your own with \emph{l'alternativa}
  \item \emph{Recall:} say the Italian for predictable, then the Italian for unpredictable, then say \emph{l'alternativa} again
\end{itemize}

\begin{tcolorbox}[breakable,colback=teal!4,colframe=teal!35!black,title={Before you move on}]
What does l'alternativa mean? (\textquotedblleft{}The alternative, the other option\textquotedblright{}.) Say it once more.
\end{tcolorbox}
\section[adattarsi]{adattarsi — to adapt, to get used to}

\label{lesson:IT-C374-adattarsi}



\begin{quote}
Before the new one: say the Italian for unpredictable, then the Italian for the alternative.
\end{quote}

\subsection*{You'll want to know: adattarsi}
\textbf{adattarsi} — \textquotedblleft{}to adapt, to get used to\textquotedblright{}.

\begin{quote}
Se cambiano i piani, mi adatto.
\end{quote}

\begin{quote}
\emph{If the plans change, I will adapt.}
\end{quote}

\emph{Adattarsi a} qualcosa. Without \emph{-si}, \emph{adattare} is to adapt something.

\begin{cousinweb}
\emph{adattarsi} is built on \emph{adattare}, a learned borrowing from Latin \emph{adaptāre} \textquotedblleft{}to fit to\textquotedblright{} (Lewis \& Short), built from \emph{ad} \textquotedblleft{}to\textquotedblright{} and \emph{aptus} \textquotedblleft{}fitted\textquotedblright{}.
\end{cousinweb}

\begin{grammarlens}[title={What you've built}]
One more word for telling, explaining and describing at B1.
\end{grammarlens}

\subsection*{Your turn}
\begin{itemize}
\raggedright
  \item \emph{Say it:} \emph{adattarsi}
  \item \emph{Say it:} \emph{adattarsi}, once more
  \item \emph{Say it:} a sentence of your own with \emph{adattarsi}
  \item \emph{Recall:} say the Italian for unpredictable, then the Italian for the alternative, then say \emph{adattarsi} again
\end{itemize}

\begin{tcolorbox}[breakable,colback=teal!4,colframe=teal!35!black,title={Before you move on}]
What does adattarsi mean? (\textquotedblleft{}To adapt, to get used to\textquotedblright{}.) Say it once more.
\end{tcolorbox}
